We will deduce this theorem from the following more general result, which we can understand as a generalization of the calculation described in Example~\ref{ex:fund-calc}:

\begin{Theorem}\label{thm:main-general}
	Let $X$ be a smooth variety or DM stack, $D\xhookrightarrow{k} X$ a divisor, and $H\xhookrightarrow{i}D$ a~smooth subvariety of codimension $2$ in $X$. Let $\frX = \Bl_HX\setminus \widetilde{D}$ and let $\cC$ denote the
homotopy pushout of dg categories
\begin{equation}\label{eq:pushout}
\colim\bigl(\Coh(H)\stackrel{i^*}\longleftarrow\Coh(D)\stackrel{k_*}\longrightarrow\Coh(X)\bigr).
\end{equation}
Then there is an equivalence
\[
\cC\cong\Coh(\frX)
\] between $\cC$ and the category of coherent sheaves on $\frX$.
\end{Theorem}
\begin{proof}
We begin by writing down the functor $\cC\to \Coh(\frX)$. From the pushout presentation of $\cC$, such a functor may be specified by the data of functors $F_H\colon \Coh(H)\to \Coh(\frX)$ and $F_X\colon \Coh(X)\to\Coh(\frX)$, together with an equivalence between the pullbacks of $F_H$ and $F_X$ to $\Coh(D)$ along the maps in the pushout \eqref{eq:pushout}.

Consider the following diagram of all the spaces involved in our story, where $E$ is the exceptional divisor of the blowup, the vertical maps are the projections, and the horizontal maps are the inclusions:
\[
\xymatrix{
\frX\ar@{^{(}->}[r]^-j & \Bl_HX\ar[d]_p & &E\ar@{_{(}->}[ll]_-{\ovi}\ar[d]^q\\
& X & D\ar@{_{(}->}[l]^k & H.\ar@{_{(}->}[l]^-i
}
\]
The functors $F_H$ and $F_X$ will be defined as the compositions
\begin{gather*}
F_H = j^*\ovi_* q^*, \qquad F_X = j^*p^*.
\end{gather*}
The restriction to $D$ of $F_X$ is evidently given by $j^*p^*k_*$, and we must produce an isomorphism between this functor and the other restriction $j^*\ovi_*q^*i^*$. Note that before applying final composition with $j^*$, these two functors differ, since $p^*k_*$ takes sheaves on $D$ to their pullback to the blowup, whereas $\ovi_* q^*i^*$ takes sheaves on $D$ to their component lying along the exceptional divisor $E$. In other words, the difference between these two functors is a sheaf on the proper transform $\widetilde{D}$ of $D$. The final pullback $j^*$ will kill this component, forcing the two functors to agree.

Thus, we have a well-defined functor $F\colon\cC\to \Coh(\frX)$.
To understand this functor,
we will need to apply
Orlov's result \cite{Orl92} that $p^*\Coh(X)$ and $q^*\Coh(H)\otimes\cO_E(1)$ generate the category $\Coh(\Bl_H(X))$ of coherent sheaves on the blowup, and in fact they form a semi-orthogonal decomposion:
\begin{equation}\label{eq:semiorth}
\Coh(\Bl_H(X)) = \langle \Coh(H)\otimes \cO_E(-1), \Coh(X)\rangle.
\end{equation}
Note that after pulling back to $\frX$, we can ignore the tensor
product with $\cO_E(-1)$ in \eqref{eq:semiorth}, since deleting $\widetilde{D}$ makes $E$ affine. This makes clear why $F$ is essentially surjective: the image of the functor $F$ contains the pullbacks to $\frX$ of both components in the semi-orthogonal decomposition~\eqref{eq:semiorth} and hence contains all objects in $j^*\Coh(\Bl_H(X))$.

To see why the functor $F$ is in fact an equivalence of categories, recall that the
category $\Coh(\frX)$ is a localization of $\Coh(\Bl_H(X))$ obtained by killing the identity morphism of any object in $\Coh\big(\widetilde{D}\big)$, and observe that this is equivalent to the localization which identifies any object in the second term of~\eqref{eq:semiorth} of the form~$k_*\cF$, for some~$\cF$ in $\Coh(D)$, with the object $(iq)^*\cF\otimes \cO_E(-1)$ in the first term.

Now note that the
category obtained from the semi-orthogonal decomposition \eqref{eq:semiorth} by identifying the images of $\Coh(X)$ and $\Coh(H)$ is equivalent to the category obtained from the ortho\-go\-nal sum $\Coh(H)\oplus \Coh(X)$ by the same identification;
but $\Coh(H)\oplus \Coh(X)$ is the~co\-pro\-duct of the categories $\Coh(H)$ and $\Coh(X)$, and the identification just described is the coequalizer of the respective maps from $D$ to these categories. In other words, this presentation exhibits $\Coh(\frX)$ as the pushout \eqref{eq:pushout}, and $F$ as the canonical equivalence $\cC\cong \Coh(\frX)$.
\end{proof}
